\section{Referente Teórico}
\label{sec:referente_teorico}

% ── A. Contextualización del campo ───────────────────────────────────────────
\subsection{Contextualización del Campo: Reconocimiento Emocional en Voz}
\label{sec:rt_contexto}

El \textit{Speech Emotion Recognition} (SER) es la disciplina que estudia métodos
automáticos para inferir el estado emocional de un hablante a partir de la señal de
voz, sin depender del contenido lingüístico. Sus fundamentos se remontan a los trabajos
de la década de 1990, cuando se observó que características prosódicas elementales
como la frecuencia fundamental (F0) y la energía presentaban correlatos sistemáticos
con estados emocionales \cite{jha2022}.

A lo largo de la primera década del siglo~XXI, el SER adoptó el paradigma de
aprendizaje automático clásico: se extraían conjuntos de características acústicas
(\textit{features}) de manera manual y se entrenaban clasificadores como SVM o Random
Forest sobre dichos vectores \cite{nath2024}. Este enfoque ofrece interpretabilidad:
es posible identificar qué características acústicas específicas contribuyen a la
decisión del modelo.

A partir de 2015, la disponibilidad de grandes corpus y de recursos de cómputo
aceleró el surgimiento de arquitecturas de aprendizaje profundo para SER. En 2023,
Wagner et al.\ demostraron que los modelos basados en \textit{transformers}
preentrenados con autosupervisión ---como Wav2Vec~2.0--- logran cerrar la denominada
\textit{valence gap}: la dificultad histórica para predecir la dimensión positivo/negativo
de la emoción \cite{wagner2023}. Sin embargo, estos modelos requieren hardware de alto
rendimiento y grandes corpus anotados, restricciones que no se satisfacen en el contexto
de un proyecto de grado con hardware estándar y un corpus de llamadas reales limitado.

Por estas razones, el presente proyecto adopta el enfoque \textbf{feature-based
clásico} combinado con detección de anomalías, estrategia que ofrece un equilibrio
apropiado entre interpretabilidad, demanda computacional y rendimiento aceptable con
corpus pequeños \cite{khare2024}.

% ── B. Teorías de las emociones ───────────────────────────────────────────────
\subsection{Teorías de las Emociones}
\label{sec:rt_teorias}

\subsubsection{Modelo dimensional de Russell}

Russell (1980) propuso el \textbf{modelo circumplejo del afecto}, que representa las
emociones en un espacio bidimensional continuo definido por dos ejes ortogonales:
\textit{valencia} (positiva a negativa) y \textit{activación} (alta a baja)
\cite{russell1980circumplex}. En este modelo, la ira y el miedo ocupan el cuadrante
de alta activación / valencia negativa, mientras que la alegría y el entusiasmo se
sitúan en alta activación / valencia positiva. La tristeza corresponde a baja
activación / valencia negativa.

Este modelo es el fundamento teórico directo de la presente investigación: el sistema
detecta \textit{picos de alta activación} (desviaciones acústicas intensas respecto
al \textit{baseline}) y posteriormente clasifica su \textit{valencia} como positiva
o negativa, operacionalizando así los dos ejes del circumplejo de Russell.

\subsubsection{Teoría de las emociones básicas de Ekman}

Ekman (1992) postuló la existencia de seis emociones básicas de carácter universal:
alegría, tristeza, ira, miedo, asco y sorpresa \cite{ekman1992basic}. Aunque este
modelo categórico ha sido ampliamente utilizado para etiquetar corpus de SER ---incluyendo
RAVDESS---, su limitación principal en el contexto de este trabajo es que las fronteras
entre categorías son difusas en el habla espontánea. La clasificación binaria positivo/negativo
del presente proyecto puede entenderse como una proyección simplificada de las emociones
básicas sobre el eje de valencia del modelo de Russell.

\subsubsection{Modelo de componentes de Scherer}

Scherer (2005) conceptualiza la emoción como un proceso multi-componente que integra
al menos cinco dimensiones: la evaluación cognitiva del estímulo, los cambios
fisiológicos, las tendencias motivacionales, la expresión motora y el sentimiento
subjetivo \cite{scherer2005what}. La pertinencia de este modelo para el presente
trabajo radica en el componente de \textbf{expresión motora vocal}: los cambios
emocionales se manifiestan en la musculatura laríngea, alterando la frecuencia,
tensión y regularidad de la vibración glótica, lo que se traduce directamente en las
características acústicas que el sistema extrae.

\subsubsection{Síntesis: por qué el modelo dimensional}

De los tres modelos, el circumplejo de Russell es el más adecuado para este proyecto
por dos razones: (a)~es continuo, lo que permite cuantificar la \textit{intensidad}
del pico emocional en términos del nivel de activación; y (b)~la dimensión de valencia
es directamente operable como etiqueta binaria. Las categorías de Ekman y los
componentes de Scherer sirven como marco interpretativo para comprender por qué
ciertas características acústicas correlacionan con los estados emocionales detectados.

% ── C. Producción vocal y paralingüística ─────────────────────────────────────
\subsection{Producción Vocal y Base Paralingüística}
\label{sec:rt_produccion}

\subsubsection{Modelo fuente-filtro}

El modelo fuente-filtro de Fant (1960) es el marco físico fundamental para comprender
la producción de voz \cite{fant1960acoustic}. Según este modelo, la voz resulta de
dos etapas: (1)~la \textbf{fuente glótica}, que genera una señal periódica cuya
frecuencia fundamental es F0, determinada por la tensión y masa de los pliegues vocales;
y (2)~el \textbf{filtro del tracto vocal}, que modifica el espectro de la fuente
mediante resonancias (formantes) que dependen de la configuración articulatoria.

Este modelo predice directamente por qué las emociones afectan la señal de voz: los
estados de alta activación emocional incrementan la tensión muscular laríngea (lo
que eleva F0 y la presión subglótica), provocan irregularidades en la vibración
glótica (lo que aumenta jitter y shimmer) y modifican la configuración del tracto
vocal (lo que altera los coeficientes MFCC).

\subsubsection{Relación emoción--fisiología vocal}

La activación emocional activa el sistema nervioso autónomo simpático, lo que produce:
(a)~incremento de la tensión de los músculos cricotiroideos, elevando F0;
(b)~aumento de la presión subglótica por mayor actividad diafragmática, elevando la
intensidad; (c)~mayor irregularidad en la aducción de los pliegues vocales, traducida
en jitter y shimmer elevados. Estos mecanismos fisiológicos son la justificación
causal ---no solo correlacional--- de las características acústicas usadas en el sistema
\cite{jha2022, lian2023}.

% ── D. Características paralingüísticas ──────────────────────────────────────
\subsection{Características Paralingüísticas}
\label{sec:rt_paralinguistico}

El análisis paralingüístico estudia los elementos de la comunicación vocal más allá
del contenido lingüístico. Específicamente, analiza el tono (altitud de la voz), el
ritmo (velocidad de habla), el volumen (intensidad sonora) y la calidad vocal
(características específicas de la voz, como si sonara ronca o clara). Estas
características transmiten información sobre el estado emocional del hablante,
independientemente del texto pronunciado. La Tabla~\ref{tab:paralinguistic} resume las
principales características, su base física en el modelo fuente-filtro y sus
implicaciones emocionales \cite{jha2022, lian2023}.

\begin{table}[!t]
\renewcommand{\arraystretch}{1.3}
\caption{Características Paralingüísticas, Fundamento Físico e Implicaciones Emocionales}
\label{tab:paralinguistic}
\centering
\begin{tabular}{p{1.6cm} p{2.0cm} p{3.8cm}}
\toprule
\textbf{Característica} & \textbf{Componente} & \textbf{Implicaciones Emocionales} \\
\midrule
Pitch (F0)       & Fuente glótica (frecuencia) & Se eleva con emociones de alta activación; desciende con baja activación. \\
Intensidad       & Fuente glótica (presión)     & Aumenta con activación emocional y estrés comunicativo. \\
Tasa de habla    & Control motor               & Se acelera con ansiedad; se ralentiza con tristeza. \\
Jitter/Shimmer   & Fuente (perturbación)        & Aumentan con tensión emocional y física \cite{zhang2023}. \\
MFCCs            & Filtro del tracto vocal      & Capturan timbre para caracterización holística \cite{jha2022}. \\
\bottomrule
\end{tabular}
\end{table}

La frecuencia fundamental (F0), la intensidad, el jitter, el shimmer y la tasa de
habla constituyen correlatos prosódicos primarios, mientras que los coeficientes
MFCC aportan información espectral sobre la configuración del tracto vocal. Lian
et al.\ (2023) confirman que la combinación de features prosódicas y espectrales
supera en rendimiento al uso de cualquier subconjunto por separado \cite{lian2023}.

% ── E. Línea base emocional ───────────────────────────────────────────────────
\subsection{Concepto de Línea Base Emocional}
\label{sec:rt_baseline}

La alta variabilidad interindividual hace necesario establecer una referencia
personalizada para cada hablante. El \textit{baseline} emocional individual es el
estado emocional neutral de esa persona en particular. Los picos se detectan como
desviaciones significativas respecto a esta \textit{baseline}, no respecto a umbrales
absolutos, lo que permite adaptarse a las características únicas de cada voz
\cite{gat2022}.

Los tres métodos más estudiados son: (1)~\textit{baseline} de inicio ---primeros
2--5~minutos de la llamada---; (2)~\textit{baseline} por mediana global ---más robusta
ante valores atípicos---; y (3)~\textit{baseline} adaptativa con ventana deslizante
---actualización continua durante la llamada---. La elección entre estos métodos
constituye una de las variables independientes del diseño experimental (ver
Sección~\ref{sec:met_diseno}).

La fundamentación teórica de la individualización del \textit{baseline} reside en la
variabilidad interhablante documentada en la literatura: Gat et al.\ (2022) demuestran
que la normalización a nivel de hablante mejora significativamente el desempeño de los
sistemas de SER en habla espontánea, ya que personas con voces naturalmente más agudas
o más intensas producirían un número excesivo de falsos positivos si se aplicaran
umbrales absolutos \cite{gat2022}.

% ── F. Modelos computacionales para SER ──────────────────────────────────────
\subsection{Modelos Computacionales para el Reconocimiento de Emociones en Voz}
\label{sec:rt_modelos}

\subsubsection{Enfoques feature-based y end-to-end}

En la literatura de SER coexisten dos grandes familias de modelos. Los modelos
\textbf{feature-based} (o de características artesanales) extraen manualmente un
vector de características acústicas a partir de la señal de voz y lo pasan a un
clasificador convencional (SVM, Random Forest, redes neuronales poco profundas).
Los modelos \textbf{end-to-end} operan directamente sobre la forma de onda o el
espectrograma usando arquitecturas profundas (CNN, LSTM, \textit{transformers}) que
aprenden representaciones internas sin intervención manual \cite{khare2024}.

Los modelos end-to-end han logrado mejoras sustanciales en benchmarks controlados,
pero presentan limitaciones prácticas: (a)~requieren grandes corpus anotados para
generalizar, (b)~son computacionalmente costosos, y (c)~su interpretabilidad es baja,
lo que dificulta justificar ante supervisores no técnicos por qué el sistema marcó un
pico específico \cite{wagner2023, delope2023}.

\subsubsection{Conjuntos de características estándar}

El conjunto \textbf{eGeMAPS} (\textit{Geneva Extended Minimalistic Acoustic Parameter
Set}), propuesto por Eyben et al.\ (2016), define 88~parámetros acústicos seleccionados
mediante criterios de relevancia teórica y eficiencia computacional
\cite{eyben2016gemaps}. Incluye estadísticas de F0, energía, formantes, jitter, shimmer
y razón armónico-ruido (HNR). Su diseño minimizaba la redundancia y maximizaba la
cobertura de los correlatos emocionales conocidos, convirtiéndolo en el estándar de
facto para evaluaciones comparativas en SER.

\subsubsection{Justificación del enfoque adoptado}

Para el presente proyecto se adopta el enfoque \textbf{feature-based} por tres
razones: (a)~\textit{Interpretabilidad clínica}: el análisis de importancia de
características (Random Forest) permite identificar qué parámetros acústicos
específicos marcaron el pico, información útil para el supervisor; (b)~\textit{Restricción
computacional}: el hardware disponible (Dell Optiplex 9020) no permite entrenar ni
desplegar modelos de aprendizaje profundo de gran escala; y (c)~\textit{Tamaño del
corpus}: con un corpus de validación de 30 llamadas, los métodos feature-based ofrecen
mejor generalización que los end-to-end, que tienden a sobreajustarse con datos
escasos \cite{khare2024, nath2024}.

% ── G. Machine learning para detección de anomalías ──────────────────────────
\subsection{Machine Learning para la Detección de Anomalías Emocionales}
\label{sec:rt_ml}

\begin{itemize}
    \item \textbf{Detección estadística:} Z-scores y umbrales de desviación respecto
          de la línea base emocional. Bajo en cómputo, altamente interpretable.
    \item \textbf{Métodos clásicos:} Support Vector Machines (SVM), Random Forest e
          Isolation Forest, algoritmos de aprendizaje automático utilizados para la
          detección de anomalías multivariadas \cite{wagner2023, xu2023, nath2024}.
          Isolation Forest construye árboles de aislamiento aleatorios y asigna
          puntuaciones de anomalía basadas en la longitud de la trayectoria de
          aislamiento \cite{xu2023}.
\end{itemize}

% ── H. Revisión de antecedentes ───────────────────────────────────────────────
\subsection{Revisión de Antecedentes}
\label{sec:rt_antecedentes}

\subsubsection{Contexto del Problema en los Call Centers}

Los call centers constituyen entornos de alta carga emocional en los que miles de
interacciones diarias pueden desembocar en picos críticos de ira, frustración o
ansiedad. La evaluación manual sigue siendo el estándar dominante, pero conlleva que
un supervisor dedique tiempo completo a cada grabación. Feng y Devillers (2023)
evidencian que los modelos entrenados con datos espontáneos de call centers superan
consistentemente a los entrenados con habla actuada \cite{feng2023}. El análisis
basado exclusivamente en el texto transcrito resulta insuficiente para capturar
estados emocionales encubiertos, como el sarcasmo o la frustración sutil, que se
manifiestan con claridad en características acústicas como la elevación del pitch, la
aceleración de la tasa de habla o el incremento de las variaciones de jitter y
shimmer \cite{jha2022, lian2023}.

\subsubsection{Características Acústicas para la Detección de Emociones}

La frecuencia fundamental (F0 o pitch), que mide la frecuencia más baja de la voz,
es uno de los correlatos acústicos más estudiados: incrementos sostenidos de F0 se
asocian con emociones de alta activación, como la ira o el miedo, mientras que
descensos de F0 se correlacionan con la tristeza o la calma \cite{jha2022, lian2023}.
Las métricas de calidad vocal jitter y shimmer miden las perturbaciones
microtemporales en la periodicidad del ciclo glótico; jitter se refiere a la
variabilidad en la frecuencia de la voz, y shimmer a la variabilidad en la amplitud.
Zhang et al.\ (2023) demostraron en IEEE que estas características capturan diferencias
prosódicas significativas bajo una alta carga emocional \cite{zhang2023}. Lian et al.\
(2023) confirman que jitter y shimmer, junto con la razón armónico-ruido (HNR),
capturan la aspereza y la ronquera asociadas con el habla ansiosa o de alta activación
emocional \cite{lian2023}. Los MFCC capturan el timbre vocal convirtiendo las señales
acústicas en una representación numérica de cómo el oído humano percibe el espectro
de frecuencia, y han sido ampliamente utilizados como entrada para clasificadores de
emociones \cite{jha2022, abdelhamid2022}. Wagner et al.\ (2023) demostraron en IEEE
TPAMI que arquitecturas de transformadores combinadas con eGeMAPS logran cerrar la
brecha de valencia en el reconocimiento de emociones en habla espontánea
\cite{wagner2023}.

\subsubsection{Normalización por Hablante y Línea Base Emocional}

La variabilidad interhablante vuelve ineficaces los umbrales absolutos y justifica el
uso de \textit{baselines} personalizadas. Gat et al.\ (2022) proponen en ICASSP un
método de normalización por hablante para sistemas de SER basados en
autosupervisión, demostrando que la normalización a nivel de hablante mejora
significativamente el desempeño en habla espontánea \cite{gat2022}. Los tres métodos
de \textit{baseline} más estudiados son: (1)~\textit{baseline} de inicio ---primeros
2--5~minutos de la llamada---; (2)~\textit{baseline} por mediana global ---más robusta
ante valores atípicos---; y (3)~\textit{baseline} adaptativa con ventana deslizante
---actualización continua durante la llamada---. La elección entre estos métodos
constituye uno de los objetivos específicos de esta investigación.

\subsubsection{Algoritmos de Machine Learning para la Detección de Picos}

El enfoque estadístico basado en Z-scores calcula la desviación de cada ventana
temporal respecto a la línea de base del hablante y marca como pico toda observación
que supere un umbral de desviación estándar. Xu et al.\ (2023) proponen Deep Isolation
Forest en IEEE TKDE, demostrando que extensiones neurales del iForest clásico superan
a los métodos lineales para la detección de anomalías en espacios de alta dimensión,
lo cual es directamente aplicable al espacio multivariado de features acústicas
\cite{xu2023}. Nath et al.\ (2024) reportan que Random Forest ofrece un equilibrio
favorable entre la interpretabilidad y la precisión en la clasificación de la valencia
emocional \cite{nath2024}. Jha et al.\ (2022) concluyen que la combinación de
características prosódicas y espectrales es más robusta que cualquiera de ellas por
separado \cite{jha2022}. Gao et al.\ (2024) presentan en ICASSP un enfoque de
fine-tuning en dos etapas con adaptadores para SER, relevante cuando el corpus de
llamadas reales disponible es limitado \cite{gao2024}.

\subsubsection{Datasets y Validación Experimental}

RAVDESS es un corpus actuado con actores profesionales ampliamente utilizado como
benchmark. Abdelhamid et al.\ (2022) reportan en IEEE Access que un CNN+LSTM entrenado
sobre RAVDESS alcanza precisiones superiores al 85\,\% en condiciones controladas,
aunque el rendimiento cae al aplicarse a la habla espontánea sin adaptación de dominio
\cite{abdelhamid2022}. IEMOCAP ofrece conversaciones elicitadas con mayor proximidad
al habla natural; Wagner et al.\ (2023) lo usan como benchmark principal para demostrar
el cierre de la brecha de valencia \cite{wagner2023}. De Lope y Graña (2023) presentan
en Neurocomputing una revisión continua del estado del arte en SER, señalando que la
generalización a dominios espontáneos es el reto abierto más relevante \cite{delope2023}.
Grósz et al.\ (2022) demuestran en ACM MM que los embeddings de Wav2Vec~2.0 transfieren
eficazmente a tareas paralingüísticas con datos limitados \cite{grosz2022}. Khare et
al.\ (2024) concluyen en Information Fusion que la combinación de características
acústicas clásicas con representaciones aprendidas es actualmente la estrategia de
mayor rendimiento en escenarios reales \cite{khare2024}.

% ── I. Articulación con la investigación ──────────────────────────────────────
\subsection{Articulación del Referente Teórico con la Investigación}
\label{sec:rt_articulacion}

El referente teórico presentado en esta sección estructura de forma directa cada uno
de los cinco objetivos específicos de la investigación:

\begin{itemize}
    \item \textbf{Objetivo 1 (extracción de características):} se sustenta en el modelo
          fuente-filtro de Fant \cite{fant1960acoustic} y en el conjunto eGeMAPS
          \cite{eyben2016gemaps}, que determinan cuáles parámetros acústicos son
          teóricamente relevantes como correlatos emocionales.
    \item \textbf{Objetivo 2 (\textit{baseline} individualizada):} se fundamenta en la
          variabilidad interhablante documentada por Gat et al.\ \cite{gat2022} y en la
          necesidad de normalización por hablante derivada del modelo circumplejo de
          Russell \cite{russell1980circumplex}, que reconoce que el punto de referencia
          emocional varía entre individuos.
    \item \textbf{Objetivo 3 (detección de anomalías):} se apoya en la teoría
          estadística de valores atípicos y en los métodos de detección de anomalías
          multivariadas revisados en la Sección~\ref{sec:rt_ml} \cite{xu2023, nath2024}.
    \item \textbf{Objetivo 4 (clasificación de valencia):} se operacionaliza
          directamente sobre el eje de valencia del modelo dimensional de Russell
          \cite{russell1980circumplex}, reduciendo el espacio emocional continuo a una
          clasificación binaria positiva/negativa.
    \item \textbf{Objetivo 5 (validación con expertos):} se inscribe en el paradigma
          empírico-positivista del SER en habla espontánea \cite{delope2023, feng2023},
          que exige la validación de sistemas automáticos contra anotaciones humanas
          en condiciones ecológicamente válidas.
\end{itemize}
